最优方案的最大能量平衡残差为 $2.27\times10^{-13}\,\mathrm{kWh}$，日末与日初储电量之差为 0，且 $\max_t(c_td_t)=0$；储电量、充放电功率和弃光约束均通过检验。纯成本模型与加入显式充放电互斥约束的对照模型所得成本一致，说明当前正电价数据下，连续 LP 已能得到物理可执行的调度。

进一步将储能容量分别下调、上调 $10\%$，购电费为 $34547.91$ 元和 $33092.98$ 元；将充放电功率上限分别下调、上调 $10\%$，购电费为 $33827.76$ 元和 $33780.88$ 元。容量变化造成的成本波动明显大于相同幅度的功率变化，表明本算例的主要限制是可跨时段转移的能量规模。固定效率未刻画功率、温度和老化的影响\cite{grimaldi}，且未计退化成本，故结果应理解为题给运行边界下的日前购电成本最优方案。
